% ============================================================
% CHARACTER REFERENCE B — Director-Level Colleague (Uni Care, Inc.)
% Form LIC 301E
% Perspective: Cross-functional director/peer leader — collaborative,
% clinical-operations perspective on systems that support elderly care.
% ============================================================
\newdocument{Character Reference --- Director-Level Colleague}
\resetlicquestions

\licformheader

\licwriterinfo
  {\RefBName}
  {\RefBAddress}
  {\RefBCityStateZip}
  {\RefBPhone}

% ----------------------------------------------------------
\licquestion{How long have you known the person you are writing this reference for?}

I have known \ApplicantName\ for \RefBKnownDuration, since \RefBKnownSince, in our shared role as senior staff members of \EmployerName. During this period, I have collaborated with him regularly on operational, clinical-technology, and compliance initiatives, giving me a well-developed understanding of his professional character.

% ----------------------------------------------------------
\licquestion{How do you know this person?}

\ApplicantName\ and I are colleagues at \EmployerName, the parent organization of \EmployerDivisions. I serve as \RefBTitle, and \ApplicantName\ serves as the \EmployerRole. Our roles intersect constantly, because the clinical services my department delivers to elderly home health and hospice patients depend on the technology infrastructure that \ApplicantName\ designs, maintains, and secures. We collaborate on the deployment of electronic health record platforms, point-of-care documentation tools used in patients' homes, HIPAA compliance initiatives, telehealth systems, staff training infrastructure, incident-response protocols, and the operational continuity of services that elderly patients rely on daily. Over the course of our working relationship, I have seen \ApplicantName\ operate under pressure, across departments, and in moments when the quality of his judgment directly affected patient-facing operations.

% ----------------------------------------------------------
\licquestion{Please give your opinion of this person's character.}

\ApplicantName\ is one of the most principled, thoughtful, and operationally disciplined professionals I have worked alongside. He approaches his responsibilities with precision, humility, and an orientation toward the people our organization ultimately serves. He is honest in his communication, consistent in his follow-through, and collaborative in how he works with departments outside his own. In a field where technology is often treated as a back-office function, he has always understood that the systems he manages are, in effect, an extension of patient care---and he carries that responsibility seriously.

I am aware of \ApplicantName's two misdemeanor DUI convictions, one from \IncidentOneDate\ and one from \IncidentTwoDate. He told me about them himself, without prompting, in a direct and honest conversation at a point when he was not required to disclose anything to me. He took full responsibility, described what he was doing to address the underlying issues, and did not attempt to soften or reframe the facts. That conversation is one of the reasons I trust him: very few people respond to personal failure with that level of candor and accountability.

What I have observed since is consistent with what he told me. He did not limit himself to the court-mandated programs. He voluntarily enrolled at \RehabProgram\ for \RehabDuration\ of intensive rehabilitation---a step no authority required of him. He has continued attending AA meetings, has made durable changes to his lifestyle and habits, and has never once permitted his personal difficulties to compromise his professional obligations. The person I work with today is a noticeably more self-aware, more emotionally regulated, and more purpose-driven version of the colleague I first met. His rehabilitation is not an abstract claim; it is visible, consistent, and sustained.

I vouch for his character without hesitation.

% ----------------------------------------------------------
\licquestion{Please describe any interaction you have observed between this person and the client group he/she is requesting to work with. For example: Clients may be children, developmentally disabled children or adults, mentally impaired adults, or elderly.}

Because my department is directly responsible for the clinical care of elderly patients in their homes, I have had many opportunities to observe how \ApplicantName\ engages with this population and with the systems that serve them:

\begin{itemize}[leftmargin=2em, itemsep=0.2em]
  \item When we have designed and rolled out clinical systems---such as EHR workflows, caregiver-family communication tools, and medication reconciliation modules---\ApplicantName\ has consistently brought the conversation back to the elderly patient: Is this workflow respectful of their time and energy? Does it preserve their dignity? Does it keep the caregiver present with the patient rather than buried in a screen? These are not questions a typical IT leader raises, and they shape the systems we deploy.
  \item I have been present with \ApplicantName\ during visits to patients' homes when on-site technology support was required. I have watched him step into those environments with quiet respect for the patient, speak softly, introduce himself to the patient and any family members present, and work efficiently to minimize disruption. On more than one occasion, I have seen him take a moment to simply sit and speak with an elderly patient who appeared to want company, before resuming his work.
  \item In clinical case conferences where technology is part of the discussion, \ApplicantName\ listens carefully before contributing. He asks informed questions about the patient's condition, family circumstances, and caregiver workflow. His participation is grounded in understanding the patient's situation, not in defending his own department.
  \item In crisis moments---a system outage during a hospice patient's final hours, a failed monitoring device at a frail patient's bedside---I have seen \ApplicantName\ respond with urgency and visible emotional seriousness, not because of a service-level agreement, but because he recognizes what is actually at stake for the person on the other end of the system.
\end{itemize}

His application for \FacilityTypeShort\ Administrator Certification is consistent with every interaction I have observed.

% ----------------------------------------------------------
\licquestion{Please add any comments you feel are relevant about this person and his/her desire to work in a community care facility.}

I want to share the following observations for the Department's consideration:

\textbf{Transparency and Self-Accountability:} \ApplicantName\ has been open with me about his record and about his rehabilitation journey. This is not a person who compartmentalizes his life or hides his history from his professional network. He has made disclosure and accountability a core operating principle, and I consider this one of the clearest indicators that his rehabilitation is genuine and enduring.

\textbf{Operational Fit for \FacilityTypeShort\ Administration:} \FacilityTypeShort\ administration is fundamentally a systems and compliance discipline executed in a human-centered environment. It requires the ability to manage complex regulatory frameworks (Title 22, CCLD oversight, medication management, incident reporting), to coordinate multidisciplinary staff, to maintain rigorous documentation under audit, and to do all of this while keeping the dignity and safety of elderly residents at the center. \ApplicantName\ has demonstrated all of these capabilities in his current role. His completion of the \RCFETraining\ has given him the regulatory-specific foundation to apply that discipline to a residential setting.

\textbf{Visible, Sustained Rehabilitation:} In the years since \IncidentTwoDate, I have observed a consistent, compounding pattern of growth in \ApplicantName---not a single dramatic change followed by stagnation, but an ongoing, daily practice of accountability, self-awareness, and healthier decision-making. This is the kind of rehabilitation that actually holds under pressure.

\textbf{A Colleague Worth Vouching For:} I understand what is at stake when the Department issues a criminal record exemption, and I would not lend my professional name to a recommendation I did not genuinely believe. Having worked closely with \ApplicantName\ for \RefBKnownDuration, with full knowledge of his record and his rehabilitation, I recommend him without reservation.

\licsignatureblock
